\chapter{What is the glue?}\label{ch:continuity}

The book's continuity is the preservation of one comparison through several mathematical languages. This chapter gathers the results into a usable hierarchy and explains the capability each adds to EBU: complete history, local computation, joint valuation, time-aware response and, where applicable, a physical probability and entropy interpretation.


\section{One comparison carried through several languages}
We began with an ordinary concern: how can an account of action say something about the physical conditions that keep life and service possible? The answer developed here is deliberately specific. Declare a state, a fixed potential and a complete physical transition. The before-minus-after difference is its EBU. Every later mathematical layer either represents that same difference or supplies a declared way to construct its endpoints.

The glue is therefore not a metaphor about balance. It is the preservation of one mathematical object across changes of representation. The chain rule connects endpoints to a path integral. Telescoping connects successive transitions. Möbius inversion connects group values to irreducible finite interactions. Smooth response and potential derivatives explain those interactions when their hypotheses hold. A canonical density can express the same values as log ratios, and the accepted thermal benchmark can express them as medium entropy changes.

Each arrow has conditions. Keeping those conditions visible is what makes the continuity rigorous rather than merely persuasive.

\diagram{continuity_chain}{The continuity hierarchy. The generator is an optional interface supplying endpoints. The main finite theory needs neither equilibrium nor an actor account. Canonical probability and P4 entropy are additional conditional branches of the same declared comparison.}{fig:continuity}

\section{The central continuity statement}
\begin{theorem}{Fixed-field EBU continuity hierarchy}\label{thm:central}
For a complete represented transition $x_0\to x_1$ and one finite, fixed potential $V$, the defining account is $E=V(x_0)-V(x_1)$. Under the additional assumptions stated beside each expression below, this same number has the following representations.
\end{theorem}

\begin{lawbox}{Central continuity: preserve the assumptions}
\textbf{D: definition.} One fixed field and complete endpoints:
\[
 E=V(x_0)-V(x_1).
\]
\textbf{P: regular path.} Theorem~\ref{thm:path}:
\[
 E=-\int_\gamma dV=-\int_\gamma\nabla V\cdot dx.
\]
\textbf{Q: canonical density.} Fixed measure, positive endpoints, $p=Z^{-1}e^{-V}$:
\[
 E=\ln[p(x_1)/p(x_0)].
\]
\textbf{T: accepted P4 scope.} Fixed conservative overdamped single-reservoir equilibrium, same temperature and no omitted work or entropy channels:
\[
 E=\Delta s_{\mathrm{med}}/k_B,
 \qquad \Delta s_{\mathrm{sys}}=-k_BE.
\]
\end{lawbox}

The proof consists of the already established chain rule, log-density identity and conditional heat balance. It does not remove any premise when the results are placed together. In particular, D does not imply Q or T. A stock account can satisfy D and P without being a thermal canonical system.

\section{The interaction form of the same continuity}
For a nonempty action set $T$, use one complete coalition table and write $G(z)=V(x(z))$. The exact finite statement is
\begin{lawbox}{Central interaction continuity}
\textbf{B: complete Boolean table.}
\[
 m_E(T)=\Delta_TE(\varnothing)=-\Delta_TG(0).
\]
\textbf{A: additive response and $C^{|T|}$ potential.}
\[
 m_E(T)=-\int_{[0,1]^{|T|}}D^{|T|}V(x(t))[h_i:i\in T]\,dt.
\]
\textbf{Q: canonical density at every corner.}
\[
 m_E(T)=\sum_{S\subseteq T}(-1)^{|T|-|S|}\ln p(x_S).
\]
\textbf{T: accepted P4 scope for every entry.}
\[
 m_{\Delta s_{\mathrm{med}}}(T)=k_Bm_E(T).
\]
\end{lawbox}

The derivative formula is one specialization of the general composite difference. With nonlinear smooth response, integrate $\partial_T(V\circ x)$ instead and use its partition chain rule. With no smooth extension, the finite table identity still stands. With a restricted feasible poset, use the applicable incidence inversion rather than assume a missing Boolean cube.

\section{What the generator contributes}
A generator supplies $x_S=\Phi_{\chi_S}^{T_f}(x_0)$ under a declared initial state, background and horizon. The finite definition values those endpoints. A measured jump or fixed increment can supply an endpoint directly; a continuous generator is one response interface rather than a prerequisite for every receipt.

\begin{lawbox}{Dynamics and interaction in the same account}
\textbf{Response.} A well-posed generator supplies the complete endpoint for each coalition.

\textbf{Temporal behaviour.} Pending state can produce memory; eigenvalues describe decay and oscillation of the declared dynamics.

\textbf{Affine ceiling.} A common linear operator with additive action inputs gives $x_S=x_\varnothing+\sum_{i\in S}h_i$. With fixed quadratic $V$, all $m_E(T)$ for $|T|\geq3$ vanish at every admissible horizon.

\textbf{Nonlinear response.} In general the exact object is $-\Delta_T(V\circ x)$. Nonlinearity can create higher orders or cancel under composition.
\end{lawbox}

These results give a service system two coordinated descriptions. Dynamics explains whether a delivery arrives, overshoots or settles. The interaction table explains how combined interventions compare with their parts. A period is not an EBU denomination, and the order of a temporal mode is not a coalition order.

\section{What is conserved, and what merely closes}
Carrier conservation concerns physical quantities within a boundary. EBU closure concerns differences of a state function. Probability normalization concerns total probability. Institutional balance concerns a declared update rule. Each identity is exact in its own setting.

For sequential receipts, the connection is
\begin{equation}
 \sum_{r=1}^{m}\bigl[V(x_{r-1})-V(x_r)\bigr]=V(x_0)-V(x_m).
\end{equation}
Intermediate values cancel because the states join. A cycle has zero total; subdividing an event preserves its total. With external changes or field revisions, their explicit lines complete the same reconciliation.

The practical consequence is reliable composition. A local receipt can join a route, a route can join a service history, and that history can be reconciled to its complete endpoint change. Joint events use their group total once. Autonomous relaxation remains in the physical history without becoming a stream of fresh actor claims.

Each discrepancy then has a place: a missing litre belongs to the carrier account, a stale baseline to the comparison, a scale revision to the field history, and an allocation dispute to the institution. This makes a large system easier to inspect and improve.

\section{Where each connection applies}
The finite comparison and Möbius reconstruction require complete declared values. The path and derivative representations add regularity. Degree ceilings add structure in both potential and response. The canonical branch adds a probability density in a fixed measure, and P4 adds the specified conservative thermal balance.

The distinction is particularly clear in the feedback tank. Its equations establish its operating point, stability, memory and affine endpoint response. They do not supply a temperature or invariant density, so its finite EBU and interaction theory stand without invoking the canonical branch. Likewise, a canonical density can be preserved by a current-carrying process, so density and reversibility remain separate properties.

This organisation lets EBU use every result at full strength within its domain. A theorem is not weakened by stating its hypotheses; the hypotheses identify the systems for which its conclusion is guaranteed.

\section{What the continuity does offer}
The completed theory gives EBU a common language for physical action at several scales. Exact cancellation makes local computation possible. Telescoping preserves the signed history. Möbius inversion explains shared dependence without adding value twice. Feedback theory includes pending commitments and timing. Degree bounds simplify the account where structure permits it. Canonical normalization connects appropriate physical models to probability and entropy.

These capabilities address concrete sources of inefficiency: duplicate claims on the same improvement, repeated dispatch against a pending delivery, omitted shared constraints, unnecessary high-order calculation and hidden changes in the yardstick. Removing those errors leaves more room for useful maintenance and more legible decisions about collective infrastructure.

The remaining chapters use this theoretical foundation to describe attribution, routes and institutional interfaces. They explain how a physical account can support cooperation, access and continuing service while keeping the choices about rights and allocation explicit. Detailed test designs, simulations and results are developed in the later books.
